\documentclass[10pt,a4paper]{amsart}
\usepackage[margin=19mm]{geometry}
\usepackage{amsmath,amssymb,mathtools}
\usepackage[scr=rsfs,cal=euler]{mathalfa}
\usepackage{xfrac}
\usepackage[unicode,hidelinks,bookmarks=false]{hyperref}
\newcommand{\ie}{i.e.\ }
\newcommand{\cf}{cf.\ }
\newcommand{\resp}{respectively\ }
\newcommand{\Subsection}{\subsection}
\providecommand{\nobreakdash}{\nobreak\hbox{-}}
\setlength{\emergencystretch}{2em}
\multlinegap0pt
\usepackage{bm}
\usepackage{bbm}
\usepackage{dsfont}
\usepackage{xspace}
\usepackage{cancel}
\usepackage{mathtools}
\usepackage{amsthm}
\makeatletter
\@ifundefined{teo}{\newtheorem{teo}{Teo}}{}
\@ifundefined{lemma}{\newtheorem{lemma}[teo]{Lemma}}{}
\makeatother
\def\norm#1{\left\|#1\right\|}
\newcommand{\ab}[1]{\left| #1 \right|}
\newcommand{\f}[1]{\mathcal{#1}}
\newcommand{\n}[1]{\mathbb{#1}}
\newcommand{\pa}[1]{\left( #1 \right)}
\title[Higher Abelian Quantum Double Models]{Higher Abelian Quantum Double Models}
\author{Jorge Acu\~na Flores, Giuseppe De Nittis, Javier Lorca Espiro}
\date{Source: arXiv:2509.12864v2 (CC BY 4.0)}
\begin{document}
\maketitle

\noindent\emph{Source and licence.} Higher Abelian Quantum Double Models. Version arXiv:2509.12864v2 (CC BY 4.0), distributed under the Creative Commons Attribution 4.0 International licence, as recorded in the arXiv licence metadata for this exact version and shown on the abstract page https://arxiv.org/abs/2509.12864v2. Full rights notice: licenses/arxiv-2509.12864-CC-BY-4.0.txt.\par\smallskip

\noindent\textit{Editorial scope.} Counted unit: the Lemma whose source label is \texttt{lem:linear\_independence\_logical} in the article (source0.tex lines 1906--1911, source label \texttt{lem:linear\_independence\_logical}), delivered together with the article's own complete original proof of it (source0.tex lines 1913--2086, one \texttt{proof} environment of 174 lines). No mathematics is written, repaired or re-derived: statement and proof are the article's own text line for line. Required context retained uncounted: lemm:lin\_indip\_g (lines 753--759); lemma:commutant\_loc (lines 1524--1531); Lem:span-K' (lines 1577--1582). Every definition, display and label of those lines is retained unchanged. Omitted from this package and not redistributed: 1--752, 760--1523, 1532--1576, 1583--1905, 1912--1912, 2087--3137. Nothing inside a retained excerpt is omitted or altered: every retained body is delivered exactly as the article's lines are written, so no omission receipt of any kind accompanies this package. Citations: the retained excerpts carry no citation command at all, so no bibliography entry of the article is transcribed and no reference is redistributed. Reference adaptation: none was needed and none was made; every reference inside the delivered unit resolves inside it, so no unresolved reference or citation is produced. Printed slips kept literal: no printed word, symbol, hypothesis or number of the counted statement or of its proof was altered. Layout and class: the article's own class is \texttt{amsart} with packages including fontenc, inputenc, babel, mathpazo, eucal, mathrsfs, microtype, amsmath, amssymb, amsthm, mathtools, upgreek, enumerate, graphicx; the wrapper is the lane's pinned \texttt{amsart} editorial wrapper at 10pt on A4, which restates the theorem-like environments the retained text needs and adds the article's own macro definitions that the retained text needs (\texttt{\textbackslash ab}, \texttt{\textbackslash f}, \texttt{\textbackslash n}, \texttt{\textbackslash norm}, \texttt{\textbackslash pa}), with extra wrapper packages (bm, bbm, dsfont, xspace, cancel, mathtools). The delivered numbering is the unit's own. Assets: no retained span contains a figure environment, a graphics-inclusion command, a PDF-page inclusion, a table or a TikZ picture; no image file is copied into this package, the package declares no asset dependency and no image file is redistributed with it. Licence: version arXiv:2509.12864v2 of this article is distributed under the Creative Commons Attribution 4.0 International licence, as recorded in the arXiv licence metadata for this exact version and shown on the abstract page https://arxiv.org/abs/2509.12864v2. A full rights notice is delivered at licenses/arxiv-2509.12864-CC-BY-4.0.txt. Characters: every retained excerpt is pure ASCII, so the delivered engine, XeLaTeX with its default Latin Modern text font, reports no missing character.
	\begin{lemma}[Linear independent generators]\label{lemm:lin_indip_g}
		Every $R\in \mathfrak{A}_{\rm loc}$ can be expressed in a unique way as finite sum
		\[
		R\;=\;\sum_{i=1}^n\lambda_{i}P_{t_i}Q_{\gamma_i}
		\]
		for some $n\in\n{N}$ with $t_i\in {\rm Hom}_0(\f{C},\n{G})^0$, $\gamma_i\in {\rm Hom}_0(\f{C},\n{G})_0$ and $\lambda_i\in\n{C}\setminus\{0\}$ for every $i=1,\ldots,n$.
	\end{lemma}

	\begin{lemma}\label{lemma:commutant_loc}
		It holds true that
		\[
		\mathfrak{K}' 
		\;=\;  \overline{\mathfrak{A}_{\rm loc}\cap\mathfrak{K}'}^{\;\|\;\|}
		\;.
		\]
	\end{lemma}

	\begin{lemma}\label{Lem:span-K'}
		It holds true that
		\[
		\mathfrak{K}'\;=\;\overline{\mathrm{span}(\mathfrak{Z})}^{\;\|\;\|}\;.
		\]
	\end{lemma}

	\begin{lemma}\label{lem:linear_independence_logical}
		The set
		$\left\{W_{\bf h}\;\middle|\;{\bf h}\in\n{H}(\f{C},\n{G})\right\}
		$
		is linearly independent, and its linear span is norm dense in $\mathfrak{A}_{\rm log}$.
	\end{lemma}

	\begin{proof}
		Consider the Hilbert space
		$
		\ell^2\pa{\n{H}(\f{C},\n{G})}
		$
		with its canonical orthonormal basis
		$
		\left\{\delta_{\bf k}\;\middle|\;{\bf k}\in\n{H}(\f{C},\n{G})\right\}.
		$
		For every ${\bf h}\in\n{H}(\f{C},\n{G})$, define the operator $V_{\bf h}$ by
		\begin{equation}\label{eq:regular_representation_H}
			V_{\bf h}\delta_{\bf k}\;:=\;\langle{\bf h},{\bf k}\rangle \delta_{{\bf h}+{\bf k}}.
		\end{equation}
		Since $\ab{\langle{\bf h},{\bf k}\rangle}=1$ and the map ${\bf k}\mapsto{\bf h}+{\bf k}$ is a bijection, $V_{\bf h}$ is a unitary operator.
		Moreover, using the properties of $\langle\,\cdot\,,\,\cdot\,\rangle$, for every ${\bf h},{\bf h}',{\bf k}\in\n{H}(\f{C},\n{G})$ one has
		\begin{align*}
			V_{\bf h}V_{{\bf h}'}\delta_{\bf k}\;=\;\langle{\bf h},{\bf h}'\rangle\langle{\bf h}',{\bf k}\rangle\langle{\bf h},{\bf k}\rangle\delta_{{\bf h}+{\bf h}'+{\bf k}}\;=\;\langle{\bf h},{\bf h}'\rangle V_{{\bf h}+{\bf h}'}\delta_{\bf k}\;.
		\end{align*}
		Hence
		\begin{equation}\label{eq:regular_Weyl_relation}
			V_{\bf h}V_{{\bf h}'}\;=\;\langle{\bf h},{\bf h}'\rangle V_{{\bf h}+{\bf h}'}\;.
		\end{equation}
		We now use these operators to construct a representation of $\mathfrak{K}'$. Let $\Lambda\subset\f{K}$ be finite and set
		$
		\mathfrak{K}'_\Lambda:=\mathfrak{A}_\Lambda\cap\mathfrak{K}'.
		$
		By Lemma~\ref{lemm:lin_indip_g} and the proof of Lemma~\ref{Lem:span-K'}, every $R\in\mathfrak{K}'_\Lambda$ admits a unique expression of the form
		$
		R=\sum_{j=1}^n c_j P_{t_j}Q_{\gamma_j}
		$,
		where
		$
		t_j\in{\rm Ker}(\delta^0)$, $\gamma_j\in{\rm Ker}(\delta_0),
		$
		and the operators $P_{t_j}Q_{\gamma_j}$ have support contained in $\Lambda$.
		We can therefore define a linear map
		\[
		\rho_\Lambda\;:\; \mathfrak{K}'_\Lambda\;\longrightarrow\;\mathcal{B}\pa{\ell^2\pa{\n{H}(\f{C},\n{G})}}
		\]
		by setting
		\begin{equation}\label{eq:rho_local_definition}
			\rho_\Lambda(P_tQ_\gamma)\;:=\;V_{([t],[\gamma])}
		\end{equation}
		on the generators and extending by linearity. The uniqueness of the above expansion shows that $\rho_\Lambda$ is well-defined.
		Let
		$
		{\bf h}=([t],[\gamma])$, ${\bf h}'=([t'],[\gamma'])
		$.
		Using the Weyl relations in $\mathfrak{K}'$, together with \eqref{eq:regular_Weyl_relation}, one obtains
		\begin{align*}
			\rho_\Lambda\pa{P_tQ_\gamma P_{t'}Q_{\gamma'}}\;&=\;\rho_\Lambda\pa{\gamma(t')P_{t+t'}Q_{\gamma+\gamma'}}\;=\;\langle{\bf h},{\bf h}'\rangle V_{{\bf h}+{\bf h}'}\\
			&=\;V_{\bf h}V_{{\bf h}'}\;
			=\;\rho_\Lambda(P_tQ_\gamma)\rho_\Lambda(P_{t'}Q_{\gamma'})\;.
		\end{align*}
		The adjoint relation is preserved as well. Indeed,
		\[
		(P_tQ_\gamma)^*\;=\;Q_{-\gamma}P_{-t}\;=\;\gamma(t)P_{-t}Q_{-\gamma}\;,
		\]
		while \eqref{eq:regular_Weyl_relation} gives
		$
		V_{\bf h}^*=\gamma(t)V_{-{\bf h}}
		$.
		Therefore
		$
		\rho_\Lambda((P_tQ_\gamma)^*)=\rho_\Lambda(P_tQ_\gamma)^*
		$.
		It follows that $\rho_\Lambda$ is a unital $\ast$-homomorphism.
		Since $\mathfrak{K}'_\Lambda$ is a $C^*$-algebra, $\rho_\Lambda$ is contractive. Thus
		\begin{equation}\label{eq:rho_local_contractive}
			\norm{\rho_\Lambda(R)}\;\leqslant\;\norm{R},\qquad R\in\mathfrak{K}'_\Lambda\;.
		\end{equation}
		The maps $\rho_\Lambda$ are compatible as $\Lambda$ increases. Indeed, if $\Lambda\subset\Lambda'$ and $R\in\mathfrak{K}'_\Lambda$, the uniqueness in Lemma~\ref{lemm:lin_indip_g} implies that the expansion of $R$ in terms of the operators $P_tQ_\gamma$ is the same whether it is regarded as an element of $\mathfrak{A}_\Lambda$ or of $\mathfrak{A}_{\Lambda'}$. Hence
		$
		\rho_{\Lambda'}(R)=\rho_\Lambda(R)
		$.
		Consequently, the family $\{\rho_\Lambda\}_\Lambda$ defines a $\ast$-homomorphism
		\[
		\rho_0\;:\;\mathfrak{A}_{\rm loc}\cap\mathfrak{K}'\;\longrightarrow\;\mathcal{B}\pa{\ell^2\pa{\n{H}(\f{C},\n{G})}}
		\]
		satisfying
		$
		\rho_0(P_tQ_\gamma)=V_{([t],[\gamma])}
		$.
		Moreover, \eqref{eq:rho_local_contractive} shows that
		$
		\norm{\rho_0(R)}\leqslant\norm{R}
		$
		for every $R\in\mathfrak{A}_{\rm loc}\cap\mathfrak{K}'$.
		By Lemma~\ref{lemma:commutant_loc}, $\mathfrak{A}_{\rm loc}\cap\mathfrak{K}'$ is norm dense in $\mathfrak{K}'$. Therefore, $\rho_0$ extends uniquely to a $\ast$-representation
		\begin{equation}\label{eq:representation_Kprime}
			\rho:\mathfrak{K}'\longrightarrow\mathcal{B}\pa{\ell^2\pa{\n{H}(\f{C},\n{G})}},
		\end{equation}
		which satisfies
		$
		\rho(P_tQ_\gamma)=V_{([t],[\gamma])}
		$.
		We now show that the representation factors through the redundancy ideal. Recall that $\mathfrak{J}$ is the norm closure of the linear span of the elements
		$
		P_tQ_\gamma-P_{t'}Q_{\gamma'}
		$
		with
		$
		t\sim t'$ and $\gamma\sim\gamma'
		$.
		The latter condition implies
		$
		([t],[\gamma])=([t'],[\gamma'])
		$,
		and therefore
		\begin{align*}
			\rho\pa{P_tQ_\gamma-P_{t'}Q_{\gamma'}}
			\;=\;V_{([t],[\gamma])}-V_{([t'],[\gamma'])}\;=\;0\;.
		\end{align*}
		By linearity and continuity of $\rho$, it follows that
		$
		\mathfrak{J}\subseteq{\rm Ker}(\rho)
		$.
		Let
		$
		\pi:\mathfrak{K}'\to\mathfrak{A}_{\rm log}=\mathfrak{K}'/\mathfrak{J}
		$
		be the canonical projection. Since $\mathfrak{J}\subseteq{\rm Ker}(\rho)$, the universal property of the quotient provides a unique $\ast$-representation
		\[
		\overline{\rho}:\mathfrak{A}_{\rm log}\longrightarrow\mathcal{B}\pa{\ell^2\pa{\n{H}(\f{C},\n{G})}}
		\]
		such that
		$
		\overline{\rho}\circ\pi=\rho
		$.
		In particular, for
		$
		{\bf h}=([t],[\gamma])
		$
		one has
		\begin{equation}\label{eq:logical_regular_representation}
			\overline{\rho}(W_{\bf h})\;=\;\overline{\rho}\pa{P_{[t]}Q_{[\gamma]}}\;=\;V_{\bf h}\;.
		\end{equation}
		Finally, suppose that for some finite subset $F\subset\n{H}(\f{C},\n{G})$ one has
		$
		\sum_{{\bf h}\in F}c_{\bf h}W_{\bf h}=0
		$.
		Applying $\overline{\rho}$ and then evaluating the resulting operator at $\delta_{\bf 0}$ gives
		\begin{align*}
			0\;&=\;\sum_{{\bf h}\in F}c_{\bf h}V_{\bf h}\delta_{\bf 0}\;
			=\;\sum_{{\bf h}\in F}c_{\bf h}\langle{\bf h},{\bf 0}\rangle\delta_{\bf h}\;
			=\;\sum_{{\bf h}\in F}c_{\bf h}\delta_{\bf h}\;.
		\end{align*}
		Since the vectors $\{\delta_{\bf h}\}_{{\bf h}\in F}$ are linearly independent, it follows that $c_{\bf h}=0$ for every ${\bf h}\in F$.
		Therefore, the operators
		$
		\left\{W_{\bf h}\;\middle|\;{\bf h}\in\n{H}(\f{C},\n{G})\right\}
		$
		are linearly independent.
		Let us now prove that the linear span of the operators $W_{\bf h}$ is norm dense in $\mathfrak{A}_{\rm log}$.
		By Lemma~\ref{Lem:span-K'}, one has
		\[
		\mathfrak{K}'=\overline{{\rm span}\left\{P_tQ_\gamma\;\middle|\;t\in{\rm Ker}(\delta^0),\ \gamma\in{\rm Ker}(\delta_0)\right\}}^{\|\cdot\|}\;.
		\]
		Since the canonical projection
		$
		\pi:\mathfrak{K}'\longrightarrow\mathfrak{A}_{\rm log}
		$
		is continuous and surjective, it follows that
		\begin{align*}
			\mathfrak{A}_{\rm log}&\;=\;\pi(\mathfrak{K}')\\
			&\subseteq\;\overline{\pi\pa{{\rm span}\left\{P_tQ_\gamma\;\middle|\;t\in{\rm Ker}(\delta^0),\ \gamma\in{\rm Ker}(\delta_0)\right\}}}^{\|\cdot\|}\\
			&=\;\overline{{\rm span}\left\{W_{\bf h}\;\middle|\;{\bf h}\in\n{H}(\f{C},\n{G})\right\}}^{\|\cdot\|}\;.
		\end{align*}
		The opposite inclusion is immediate. Therefore,
		\[
		\mathfrak{A}_{\rm log}\;=\;\overline{{\rm span}\left\{W_{\bf h}\;\middle|\;{\bf h}\in\n{H}(\f{C},\n{G})\right\}}^{\|\cdot\|}
		\]
		and this concludes the proof.	
	\end{proof}

\end{document}
